\documentclass[../../main.tex]{subfiles}% 注意这里的文件路径不能用 ./main.tex ,否则用latexmk编译子文件会报错
\graphicspath{{\subfix{./image/}}{\subfix{../../image/}}} % 指定图片引用路径,后续可以直接使用图片文件名
% 如果图片存放在上级目录,则使用 ../../image/ 作为备用路径

% 例如:
% \begin{figure}[H]
% \centering
% \includegraphics[width=0.8\textwidth]{图.png}
% \caption{}
% \label{figure:图}
% \end{figure}
% 注意:上述\label{}一定要放在\caption{}之后,否则引用图片序号会只会显示??.

\begin{document}

\section{重积分不等式证明}

\begin{example}
\begin{enumerate}[(1)]
\item 设$u\in C^1[0,1]$满足$u(0)=0$，证明：
\begin{align*}
|u(1)|^2 \leqslant \tanh 1\cdot \left( \int_{0}^{1}\bigl[|u(r)|^2+|u'(r)|^2\bigr]\mathrm{d}r \right).
\end{align*}

\item 设$B\subset\mathbb{R}^3$是单位球，证明：
\begin{align*}
\int_{\partial B}|f|^2\mathrm{d}S\leqslant \frac{e^2-1}{2}\left( \int_{B}|f|^2\mathrm{d}V+\int_{B}|\nabla f|^2\mathrm{d}V \right),\quad \forall f\in C^\infty(\overline{B}),
\end{align*}
并证明系数不可改进.
\end{enumerate}
\end{example}
\begin{remark}
\begin{enumerate}[(1)]
\item 第(1)问中,为了分部积分能表示出$u(1)$，我们只需令$\eta$满足$\eta''=\eta$，$\eta'(1)=1$。

\item 第(2)问的思路是:用极坐标将对球体的积分转化为径向部分的函数的积分.
\end{enumerate}
\end{remark}
\begin{proof}\begin{enumerate}[(1)]
\item 令$\eta(r)\triangleq \dfrac{\sinh r}{\cosh 1}$，则
\begin{align*}
\eta'(r)=\frac{\cosh r}{\cosh 1},\quad \eta''(r)=\frac{\sinh r}{\cosh 1}=\eta(r).
\end{align*}
由分部积分知
\begin{align*}
\int_{0}^{1}u'(r)\eta'(r)\mathrm{d}r
=u(1)\eta'(1)-\int_{0}^{1}u(r)\eta''(r)\mathrm{d}r
=u(1)-\int_{0}^{1}u(r)\eta(r)\mathrm{d}r.
\end{align*}
从而
\begin{align*}
u(1)=\int_{0}^{1}\bigl(u'(r)\eta'(r)+u(r)\eta(r)\bigr)\mathrm{d}r.
\end{align*}
于是由Cauchy不等式和\refpro{proposition:双曲三角函数基本性质}可得
\begin{align*}
|u(1)|^2
&\leqslant \left| \int_{0}^{1}\bigl(u'(r)\eta'(r)+u(r)\eta(r)\bigr)\mathrm{d}r \right|^2 \\
&\overset{\text{\hyperref[theorem:Cauchy不等式(离散版)]{Cauchy不等式(离散版)}}}{\leqslant}
\left| \int_{0}^{1}\sqrt{|u'(r)|^2+|u(r)|^2}\cdot \sqrt{|\eta'(r)|^2+|\eta(r)|^2}\,\mathrm{d}r \right|^2 \\
&\overset{\text{\hyperref[theorem:Cauchy积分不等式]{Cauchy积分不等式}}}{\leqslant}
\left( \int_{0}^{1}\bigl(|u'(r)|^2+|u(r)|^2\bigr)\mathrm{d}r \right)
\left( \int_{0}^{1}\bigl(|\eta'(r)|^2+|\eta(r)|^2\bigr)\mathrm{d}r \right) \\
&=\left( \int_{0}^{1}\bigl(|u'(r)|^2+|u(r)|^2\bigr)\mathrm{d}r \right)
\int_{0}^{1}\frac{\cosh^2 r+\sinh^2 r}{\cosh^2 1}\mathrm{d}r \\
&=\left( \int_{0}^{1}\bigl(|u'(r)|^2+|u(r)|^2\bigr)\mathrm{d}r \right)
\int_{0}^{1}\frac{\cosh(2r)}{\cosh^2 1}\mathrm{d}r \\
&=\left( \int_{0}^{1}\bigl(|u'(r)|^2+|u(r)|^2\bigr)\mathrm{d}r \right)\cdot \frac{\sinh 2}{2\cosh^2 1} \\
&=\tanh 1\cdot \int_{0}^{1}\bigl(|u'(r)|^2+|u(r)|^2\bigr)\mathrm{d}r.
\end{align*}

\item 对$\forall w>0$，令$F_w(r)\triangleq rf(rw) \ (\forall r\in[0,1])$，则利用$(1)$可得
\begin{align*}
|f(w)|^2&=|F_w(1)|^2\leqslant \tanh 1\int_{0}^{1}\left(|F_w(r)|^2+|F_w'(r)|^2\right)\mathrm{d}r\\
&=\tanh 1\int_{0}^{1}\left(r^2|f(rw)|^2+\left|f(rw)+r\frac{\partial f(rw)}{\partial r}\right|^2\right)\mathrm{d}r\\
&=\tanh 1\int_{0}^{1}\left[r^2|f(rw)|^2+r^2\left|\frac{\partial f(rw)}{\partial r}\right|^2+|f(rw)|^2+2r^2f(rw)\frac{\partial f(rw)}{\partial r}\right]\mathrm{d}r\\
&=\tanh 1\int_{0}^{1}\left[r^2|f(rw)|^2+r^2\left|\frac{\partial f(rw)}{\partial r}\right|^2+\frac{\mathrm{d}\left(r|f(rw)|^2\right)}{\mathrm{d}r}\right]\mathrm{d}r\\
&=\tanh 1\int_{0}^{1}\left(r^2|f(rw)|^2+r^2\left|\frac{\partial f(rw)}{\partial r}\right|^2\right)\mathrm{d}r+\tanh 1|f(w)|^2.
\end{align*}
从而
\begin{align*}
|f(w)|&\leqslant \frac{\tanh 1}{1-\tanh 1}\int_{0}^{1}\left(r^2|f(rw)|^2+r^2\left|\frac{\partial f(rw)}{\partial r}\right|^2\right)\mathrm{d}r\\
&=\frac{e^2-1}{2}\int_{0}^{1}\left(r^2|f(rw)|^2+r^2\left|\frac{\partial f(rw)}{\partial r}\right|^2\right)\mathrm{d}r.
\end{align*}
于是
\begin{align*}
\int_{\partial B}|f(w)|^2\mathrm{d}S(w)&\leqslant \frac{e^2-1}{2}\int_{\partial B}\int_{0}^{1}\left(r^2|f(rw)|^2+r^2\left|\frac{\partial f(rw)}{\partial r}\right|^2\right)\mathrm{d}r\mathrm{d}S(w)\\
&=\frac{e^2-1}{2}\int_{0}^{1}\int_{|w|=1}\left(r^2|f(rw)|^2+r^2\left|w\cdot \nabla f(rw)\right|^2\right)\mathrm{d}S(w)\mathrm{d}r\\
&\xlongequal[\ v=rw\ ]{\text{\hyperref[theorem:第一类曲线曲面积分换元法]{第一类曲面积分换元法}}}\frac{e^2-1}{2}\int_{0}^{1}\int_{|v|=r}\left(|f(v)|^2+\left|\frac{v}{r}\cdot \nabla f(v)\right|^2\right)\mathrm{d}S(v)\mathrm{d}r\\
&=\frac{e^2-1}{2}\int_{0}^{1}\int_{|v|=r}\left(|f(v)|^2+\left|\nabla f(v)\right|^2\right)\mathrm{d}S(v)\mathrm{d}r\\
&\xlongequal{\text{\hyperref[theorem:余面积公式]{余面积公式}}}\frac{e^2-1}{2}\iiint_{0\leqslant|v|\leqslant 1}\left(|f(v)|^2+|\nabla f(v)|^2\right)\mathrm{d}V\\
&=\frac{e^2-1}{2}\int_B\left(|f|^2+|\nabla f|^2\right)\mathrm{d}V.
\end{align*}
特别地，当
\begin{align*}
f(x,y,z)=\frac{\sinh \sqrt{x^2+y^2+z^2}}{\sqrt{x^2+y^2+z^2}}
\end{align*}
时，不等式等号成立.
\end{enumerate}
\end{proof}













\end{document}